\documentclass[11pt,a4paper]{article}

\usepackage{authblk}
\usepackage[utf8]{inputenc}
\usepackage[T1]{fontenc}
\usepackage[spanish,es-nodecimaldot]{babel}
\usepackage{lmodern}
\usepackage{microtype}
\usepackage{amsmath,amssymb,amsthm,mathtools}
\usepackage{geometry}
\geometry{margin=2.5cm}
\usepackage{booktabs}
\usepackage{tabularx}
\usepackage{array}
\usepackage{enumitem}
\usepackage{alltt}
\usepackage{xcolor}
\usepackage{listings}
\usepackage{hyperref}
\hypersetup{
  colorlinks=true,
  linkcolor=blue,
  urlcolor=blue,
  citecolor=blue
}

\newcommand{\N}{\mathbb{N}}
\newcommand{\Z}{\mathbb{Z}}
\newcommand{\R}{\mathbb{R}}
\DeclareMathOperator*{\argmax}{arg\,max}

\lstdefinestyle{codigo}{
  basicstyle=\ttfamily\small,
  breaklines=true,
  columns=fullflexible,
  keepspaces=true,
  showstringspaces=false,
  frame=single,
  extendedchars=true,
  inputencoding=utf8
}
\lstset{style=codigo}

\title{Manual técnico: Aprendizaje Federado}
\author{Alejandro Martínez Orozco y Jon Henry Vásquez Mesa}

\affil{Análisis y Diseño de Algoritmos $\cdot$ Grupo 2}
\affil{Universidad de Caldas}

\date{6 de noviembre de 2026}

\begin{document}

\maketitle

\tableofcontents
\newpage

%=====================================================================
\section{Propósito y alcance}
%=====================================================================

La solución base resuelve las dos partes del problema para el caso de la
docente (\texttt{data/N15.json}):

\begin{itemize}[leftmargin=*,label=$\bullet$]
  \item \textbf{Parte 1, reparto del presupuesto de ruido:} dados los nodos ya
  elegidos y un presupuesto $B$, repartirlo para maximizar
  $\sum_i u_i(x_i)$.
  \item \textbf{Parte 2, selección de participantes:} elegir en cada ronda el
  subconjunto de nodos que maximiza una utilidad conjunta submodular, con un
  cupo $k$ y un presupuesto de costos de participación $C$.
\end{itemize}

Las dos partes se encadenan a lo largo de $T$ rondas. La selección decide quién
participa, el reparto decide cuánto ruido gasta cada uno y el presupuesto total
de cada nodo, $P_i$, se va gastando de una ronda a otra.

Quedan para etapas posteriores del cronograma los fallos de nodos, las
externalidades de red en el reparto, el algoritmo exacto con su $T(n)$, la cota
de aproximación con costos y la validación experimental con instancias
crecientes.

%=====================================================================
\section{Contexto e investigación previa}
%=====================================================================

\subsection{Por qué el problema existe}

En aprendizaje federado (McMahan et al., 2017) varios nodos entrenan un modelo
común sin compartir sus datos crudos, y solo envían actualizaciones. Esas
actualizaciones todavía pueden filtrar información, así que se protegen con
privacidad diferencial, que añade ruido calibrado (Dwork y Roth, 2014). Más
ruido protege más, pero reduce la utilidad de la actualización, y cada nodo
tolera solo una cantidad total de ruido (su presupuesto de privacidad). En cada
ronda hay entonces dos decisiones: quién participa y cuánto de su presupuesto
gasta.

\subsection{Qué se sabe de cada parte}

\begin{table}[htbp]
\centering
\small
\renewcommand{\arraystretch}{1.2}
\begin{tabularx}{\textwidth}{@{}>{\raggedright\arraybackslash}p{4.6cm}>{\raggedright\arraybackslash}X>{\raggedright\arraybackslash}p{3.6cm}@{}}
\toprule
Pregunta & Resultado conocido & Fuente \\
\midrule
Repartir $B$ unidades discretas entre $N$ funciones separables, cóncavas y
crecientes &
Es un problema de asignación de recursos con solución exacta en tiempo
polinomial: el voraz por ganancia marginal es óptimo &
Ibaraki y Katoh (1988) \\
\addlinespace
Elegir $k$ elementos para maximizar una función monótona y submodular &
NP-difícil; el voraz garantiza $\ge(1-1/e)\approx63{,}2\,\%$ del óptimo &
Nemhauser, Wolsey y Fisher (1978) \\
\addlinespace
¿Se puede mejorar esa garantía? &
No, salvo P $=$ NP, ni siquiera para Cobertura Máxima &
Feige (1998) \\
\addlinespace
Con costos y presupuesto (\emph{budgeted maximum coverage}) &
El voraz simple pierde la garantía; hay algoritmos que la recuperan &
Khuller, Moss y Naor (1999); Sviridenko (2004) \\
\addlinespace
Dureza de la selección &
Reducción desde Cobertura Máxima, a su vez desde Cobertura de Conjuntos &
Karp (1972); \texttt{Formalización.pdf}, Teorema de NP-dureza \\
\bottomrule
\end{tabularx}
\caption{Resultados conocidos que sustentan el diseño de la solución.}
\end{table}

\subsection{Qué decisiones tomamos a partir de eso}

\begin{itemize}[leftmargin=*,label=$\bullet$]
  \item La Parte 1 se resuelve de forma exacta con el voraz de montículo.
  Existe un algoritmo óptimo y eficiente, así que no tiene sentido recurrir a
  fuerza bruta ni a metaheurísticas.
  \item La Parte 2 usa un voraz: la selección es NP-difícil y la función es
  submodular, que es el caso en el que Nemhauser et al.\ dan garantía.
  Mientras $N$ lo permite, comparamos además con el óptimo por fuerza bruta.
  \item El caso trae costos distintos además del cupo, y con ellos la cota
  $(1-1/e)$ no aplica tal cual. Por eso implementamos dos variantes del voraz
  y damos la comparación con el óptimo solo como referencia (sección~12).
\end{itemize}

%=====================================================================
\section{Problema formalizado (resumen)}
%=====================================================================

La formalización completa (entradas, salidas, tamaño $n$, objetivo,
restricciones, casos límite) está en \texttt{doc/Formalización.pdf}. Para el
caso G2, por ronda $t$:

\begin{itemize}[leftmargin=*,label=$\bullet$]
  \item \textbf{Entrada:} $N$ nodos, cupo $k$, costos $c_i$, presupuesto de
  participación $C$, ruido $B$, presupuestos totales $P_i$, grafo de influencia
  $w_{ij}$, funciones $u_i$ y $f$.
  \item \textbf{Decisión:} $S_t\subseteq V$ y
  $x^t\in\Z_{\ge0}^{S_t}$.
  \item \textbf{Restricciones:} $|S_t|\le k$;
  $\sum_{i\in S_t} c_i\le C$; $\sum_i x_i^t\le B$;
  $\sum_t x_i^t\le P_i$.
  \item \textbf{Objetivo:} maximizar
  $\sum_t\bigl[f(S_t)+\sum_{i\in S_t}u_i(x_i^t)\bigr]$.
  \item \textbf{Tamaño de la entrada:} $n=N$ (el grafo del caso es completo,
  así que $|E|=N(N-1)/2=105$; $k$, $B$, $C$, $T$ se reportan aparte).
\end{itemize}

Para el caso: $N=15$, $T=3$, $k=5$, $B=30$, $C=24$, $P_i=18$ para todos,
costos entre 2 y 7.

%=====================================================================
\section{Decisiones de modelado y su justificación}
%=====================================================================

Se eligieron las siguientes fórmulas de utilidad:

\subsection{Utilidad de reparto: $u_i(x)=a_i\ln(1+x)$}

\begin{itemize}[leftmargin=*,label=$\bullet$]
  \item $a_i$ es \texttt{coeficiente\_utilidad}.
  \item Es creciente y cóncava, con $u_i(0)=0$, como exige el modelo.
  \item Ganancia marginal de la unidad $x+1$:
  \[
    \Delta_i(x+1)=a_i\ln\frac{x+2}{x+1},
  \]
  estrictamente decreciente en $x$. Es un caso de rendimientos decrecientes:
  las primeras unidades de presupuesto aportan más que las siguientes.
  \item Se usa \texttt{log1p(1/(x+1))} para calcularlo con estabilidad numérica.
\end{itemize}

\subsection{Utilidad conjunta: definición explícita de $f$}

Dados un conjunto de objetivos $J=\{1,\dots,m\}$, pesos $\beta_j\ge0$ y
probabilidades $w_{ij}\in[0,1]$:
\[
  f(S)=\sum_{j\in J}\beta_j\left(1-\prod_{i\in S}(1-w_{ij})\right).
\]
Es una cobertura probabilística ponderada: cada nodo seleccionado $i$
<<cubre>> al objetivo $j$ con probabilidad $w_{ij}$, y cubrir a $j$ vale
$\beta_j$. En el caso G2, $J=V$ (así $m=N$), $\beta_j=a_j$, $w_{ij}$ es la
\texttt{peso\_influencia} del grafo y $w_{jj}=1$ (un nodo seleccionado se cubre
a sí mismo). Con $w_{ij}\in\{0,1\}$ y $\beta_j=1$ se obtiene
$\bigl|\bigcup F_i\bigr|$, la Cobertura Máxima.

\emph{Por qué es monótona y submodular.} Sea
$\mu_S(j)=\prod_{i\in S}(1-w_{ij})\in[0,1]$, la probabilidad de que $j$ siga
sin cubrir. Agregar un nodo $i$ a $S$ aporta
\[
  f(S\cup\{i\})-f(S)=\sum_j\beta_j\,\mu_S(j)\,w_{ij}\ge0,
\]
lo que prueba monotonía. Si $A\subseteq B$, entonces $\mu_B(j)\le\mu_A(j)$
para todo $j$, porque se multiplican más factores en $[0,1]$. Por tanto, la
ganancia de $i$ sobre $A$ es mayor o igual que sobre $B$, que es la
submodularidad. Guardando el vector $\mu(\cdot)$, las ganancias se calculan de
forma incremental.

Esta $f$ encaja con la reducción de NP-dureza, ya que la cobertura
$\bigl|\bigcup F_i\bigr|$ es el caso $w_{ij}\in\{0,1\}$. El grafo del caso
influye en la selección a través de $f$.

\subsection{Otras decisiones}

\begin{itemize}[leftmargin=*,label=$\bullet$]
  \item La influencia es simétrica ($w_{ij}=w_{ji}$), como en el caso, que trae
  aristas no dirigidas.
  \item No modelamos fallos de nodos, porque el caso no trae probabilidades
  $q_i$. Un nodo está disponible mientras le quede presupuesto.
  \item El presupuesto es entero y discreto, y el reparto asigna las unidades
  de una en una.
  \item Cupo y costos rigen a la vez: el caso da $k=5$ y $C=24$, y se exigen
  ambos.
  \item El presupuesto $P_i$ se gasta entre rondas y no se transfiere entre
  nodos.
  \item Los empates se resuelven por menor índice de nodo, para que la
  ejecución sea determinista.
\end{itemize}

%=====================================================================
\section{Algoritmos}
%=====================================================================

\subsection{Reparto voraz con montículo (Parte 1)}

\begin{alltt}
reparto(S, B, restante):
    x[i] \(\leftarrow\) 0  para i en S
    H \(\leftarrow\) montículo máximo con \((\Delta_i(1), i)\) para i en S con restante[i] > 0
    repetir B veces:
        si H vacío: parar
        (g, i) \(\leftarrow\) extraer máximo de H
        x[i] \(\leftarrow\) x[i] + 1
        si x[i] < restante[i]: insertar \((\Delta_i(x[i]+1), i)\) en H
    devolver x
\end{alltt}

\textbf{Corrección (argumento de intercambio).} Sea $x^*$ una solución óptima
con $\sum x_i^*=B$ que respeta los topes, y supongamos que no es la del voraz.
Existen $i,j$ con $x_i^*>g_i$ y $x_j^*<g_j$, donde $g$ es la salida del voraz,
tales que la última unidad de $i$ en $x^*$ tiene ganancia $\Delta_i(x_i^*)$ y la
siguiente de $j$, $\Delta_j(x_j^*+1)$. El voraz tomó antes la unidad
$x_j^*+1$ de $j$ que la unidad $x_i^*$ de $i$, por lo que
$\Delta_j(x_j^*+1)\ge\Delta_i(x_i^*)$. Mover una unidad de $i$ a $j$ (los topes
se respetan, pues $x_j^*+1\le g_j\le\mathit{restante}[j]$) no disminuye la
utilidad total. Repitiendo, se llega a $g$ sin perder optimalidad. La concavidad
(ganancias no crecientes) garantiza que el orden de ganancias de cada nodo es el
orden en que se asignan las unidades. Los topes $P_i$ (cotas superiores enteras)
no rompen el argumento.

\subsection{Selección voraz con cupo y costos (Parte 2)}

\begin{alltt}
seleccion(elegibles, modo):
    S \(\leftarrow \emptyset\) ; costo \(\leftarrow\) 0 ; \(\mu[j] \leftarrow 1\) para todo j
    mientras \(|S| < k\):
        cand \(\leftarrow \{ i \in \mathit{elegibles} \setminus S : \mathit{costo} + c_i \le C \}\)
        si cand vacío: parar
        para i en cand: \(g_i \leftarrow \sum_j \beta_j \cdot \mu[j] \cdot w_{ij}\)   \# ganancia marginal
        elegir \(i^* = \argmax g_i\)              (modo "ganancia")
               o  \(\argmax g_i / c_i\)          (modo "razón")
        S \(\leftarrow S \cup \{i^*\}\) ; costo \(\leftarrow\) costo + \(c_{i^*}\)
        \(\mu[j] \leftarrow \mu[j] \cdot (1 - w_{i^*,j})\) para todo j
    devolver S
\end{alltt}

En cada ronda se ejecutan los dos modos y se conserva el de mayor $f(S)$.

\textbf{Qué garantiza y qué no.}

\begin{itemize}[leftmargin=*,label=$\bullet$]
  \item Sin costos ($c_i=1$, $C=k$), el voraz por ganancia cumple
  $f(S)\ge(1-1/e)\,f(S_{\mathrm{opt}})$ (Nemhauser et al., 1978). La idea es
  que, en cada paso, la ganancia marginal del voraz es al menos $1/k$ de lo que
  falta para llegar al óptimo, lo que da
  \[
    f(S_k)\ge\bigl(1-(1-1/k)^k\bigr)f(S_{\mathrm{opt}})
    \ge(1-1/e)f(S_{\mathrm{opt}}).
  \]
  \item Con costos distintos \textbf{no hay garantía}. Ejemplo: $C=10$, un nodo
  $A$ de ganancia 1 y costo 1, y un nodo $B$ de ganancia 1{,}1 y costo 10. El
  voraz por ganancia toma $B$ y agota el presupuesto, mientras que el óptimo
  toma $A$ y otros nodos baratos por mucho más; con ganancias de ese estilo la
  razón puede empeorar tanto como se quiera. El voraz por razón evita este
  caso, pero tiene su propio contraejemplo: un nodo muy barato y de poca
  ganancia le gana a uno de gran valor que sí cabía.
  \item Para \emph{budgeted maximum coverage} hay garantías que salen de dos
  correcciones al voraz por razón. La solución base no implementa ninguna, y
  quedan para la etapa del algoritmo aproximado. Se explican a continuación.
\end{itemize}

\subsubsection{Corrección 1: añadir el mejor nodo individual (Khuller, Moss y Naor, 1999)}

\emph{Qué es.} Además del voraz por razón ganancia/costo, se calcula el mejor
nodo individual: el nodo $i$ con $c_i\le C$ que maximiza $f(\{i\})$. La
respuesta es el mejor entre la solución del voraz y ese nodo solo. Cuesta $N$
evaluaciones de $f$ más, $O(N^2)$ con esta $f$.

\emph{Por qué hace falta.} El voraz por razón puede llenar el presupuesto con
nodos baratos y dejar fuera un nodo caro que valía muchísimo. Ejemplo con
$C=10$:

\begin{table}[htbp]
\centering
\small
\begin{tabular}{@{}lrrr@{}}
\toprule
Nodo & Ganancia & Costo & Razón \\
\midrule
A & 2  & 1  & 2{,}0 \\
B & 10 & 10 & 1{,}0 \\
\bottomrule
\end{tabular}
\caption{Ejemplo en el que el voraz por razón falla.}
\end{table}

El voraz por razón toma $A$ primero, que tiene la razón mayor. Después $B$ ya
no cabe ($1+10>10$), así que obtiene 2 frente a un óptimo de 10. Si $B$ costara
$C$ y valiera $C$, la razón entre lo obtenido y el óptimo sería $2/C$, tan mala
como se quiera. El mejor nodo individual arregla este caso, porque $B$ solo
vale 10. Con esta corrección se prueba una garantía de
$(1-1/e)/2\approx0{,}316$ del óptimo.

\subsubsection{Corrección 2: enumerar los tripletes iniciales (Sviridenko, 2004)}

\emph{Qué es.} Se prueba cada subconjunto de hasta 3 nodos que quepa en $C$
como semilla. Para cada semilla se completa con el voraz por razón sobre los
nodos restantes, y se conserva la mejor solución de todas. Son $O(N^3)$
semillas.

\emph{Por qué ayuda.} El óptimo contiene algunos nodos de mucho valor. Entre
todas las semillas de tamaño 3 hay una que coincide con los 3 nodos de mayor
valor del óptimo, y al voraz solo le queda completar el resto, donde sus errores
pesan menos. La garantía sube a $(1-1/e)\approx0{,}632$, igual que sin costos.

\emph{Costo.} Cada una de las $O(N^3)$ semillas lanza un voraz de
$O(k\cdot N^2)$ con esta $f$, o sea $O(k\cdot N^5)$ en total. Es polinomial
pero caro, y por eso no está en la solución base.

\subsubsection{Qué hace entonces la solución base}

Ejecuta el voraz por ganancia y el voraz por razón, y se queda con el mejor.
Eso no equivale a ninguna de las dos correcciones, porque no mira el mejor nodo
individual ni enumera semillas. Por eso no podemos afirmar ninguna de las dos
garantías. En el caso entregado el resultado coincide con el óptimo
(razón 1,0), pero es un resultado de esa instancia, no una garantía. Añadir el
mejor nodo individual es barato ($N$ evaluaciones por ronda) y es el siguiente
paso.

%=====================================================================
\section{Análisis de complejidad}
%=====================================================================

Notación: $N$ nodos, $k$ cupo, $B$ presupuesto de ruido por ronda, $T$ rondas.

\paragraph{Reparto (por ronda).} Construir el montículo: $O(|S|)$ con
$|S|\le k$. Cada una de las $B$ extracciones y reinserciones cuesta
$O(\log|S|)$. Total: $O(k+B\log k)$. Es polinomial en $B$ y $N$, y
pseudo-polinomial en la longitud de la codificación binaria de $B$ (el informe
lo aclara). En la formulación del curso, con $S=N$: $T(N,B)=O(B\log N)$.

\paragraph{Selección (por ronda).}

\begin{itemize}[leftmargin=*,label=$\bullet$]
  \item Ganancia marginal de un candidato con esta $f$: recorre los $N$ nodos
  $j$ $\to$ $O(N)$.
  \item Un paso evalúa hasta $N$ candidatos $\to$ $O(N^2)$; actualizar
  $\mu(\cdot)$ cuesta $O(N)$.
  \item $k$ pasos $\to$ $O(k\cdot N^2)$; dos modos $\to$ factor constante 2.
  \item Como $f$ está definida explícitamente, la ganancia marginal cuesta
  $O(m)$, y aquí $m=N$. En general el costo es $O(k\cdot N\cdot m)$. No usamos
  el $O(k\cdot N)$ que daría un oráculo $O(1)$, porque ya no hay oráculo, y la
  $T(n)$ experimental se contrastará con esta expresión.
\end{itemize}

\paragraph{Total.} $T$ rondas:
$O\bigl(T\cdot(k\cdot N^2+k+B\log k)\bigr)$. Con los parámetros del caso fijos
es $O(N^2)$, y el caso real se resuelve en cerca de un milisegundo.

\paragraph{Fuerza bruta de verificación.} Enumera
$\sum_{r\le k}\binom{N}{r}$ subconjuntos y evalúa $f$ en $O(r\cdot N)$ cada
uno. Para $N=15$, $k=5$ son 4\,944 subconjuntos por ronda. Se omite si hay más
de 300\,000, y \texttt{-{}-sin-exacto} la omite siempre.

%=====================================================================
\section{Estructura del código}
%=====================================================================

\begin{alltt}
src/red.py            solución base (un solo módulo, sin dependencias)
tests/test\_red.py     pruebas automáticas
data/N15.json         caso de la docente
resultados/           resultado.json y traza.txt por caso
doc/                  informe (.pdf y .tex) y este manual
\end{alltt}

Funciones de \texttt{src/red.py}:

\begin{table}[htbp]
\centering
\small
\renewcommand{\arraystretch}{1.2}
\begin{tabularx}{\textwidth}{@{}>{\raggedright\arraybackslash}p{5.6cm}>{\raggedright\arraybackslash}X@{}}
\toprule
Elemento & Responsabilidad \\
\midrule
\texttt{Instancia} &
Lee el JSON; construye la matriz de influencia (simétrica, $w_{ii}=1$);
expone \texttt{ganancia\_unidad}, \texttt{utilidad\_nodo} y \texttt{f} \\
\addlinespace
\texttt{seleccion\_voraz(inst, elegibles, modo)} &
Voraz con cupo y costos; devuelve $S$ y traza de cada paso con todos los
candidatos \\
\addlinespace
\texttt{seleccion\_exacta(inst, elegibles)} &
Fuerza bruta de verificación \\
\addlinespace
\texttt{reparto\_voraz(inst, S, B, restante)} &
Voraz de montículo respetando topes; devuelve $x$ y traza por unidad \\
\addlinespace
\texttt{resolver(inst)} &
Orquesta las $T$ rondas, actualiza presupuestos restantes y calcula la utilidad
total \\
\addlinespace
\texttt{verificar(...)} &
Comprueba restricciones y, si hay exacto, la razón al óptimo \\
\addlinespace
\texttt{texto(...)}, \texttt{main()} &
Traza legible y línea de comandos \\
\bottomrule
\end{tabularx}
\caption{Responsabilidades de las funciones principales.}
\end{table}

%=====================================================================
\section{Formatos de entrada y salida}
%=====================================================================

\subsection{Entrada (JSON)}

\begin{lstlisting}
{
  "parametros": {
    "numero_nodos": 15, "numero_rondas": 3,
    "cupo_participantes_por_ronda_k": 5,
    "presupuesto_ruido_por_ronda_B": 30,
    "presupuesto_costos_participacion_por_ronda": 24
  },
  "grafo": [ {"u": "N01", "v": "N02", "peso_influencia": 0.25}, ... ],
  "nodos": [ {"id": "N01", "coeficiente_utilidad": 8.5,
              "presupuesto_total_nodo": 18, "costo_participacion": 3}, ... ]
}
\end{lstlisting}

Todos los campos son obligatorios. El programa comprueba que
\texttt{numero\_nodos} coincida con la lista de nodos. Cualquier otro error de
formato termina con la excepción de Python correspondiente, porque no hay
validación de esquema más fina.

\subsection{Salida}

\texttt{resultados/<caso>\_resultado.json}:

\begin{table}[htbp]
\centering
\small
\renewcommand{\arraystretch}{1.2}
\begin{tabularx}{\textwidth}{@{}>{\raggedright\arraybackslash}p{6.2cm}>{\raggedright\arraybackslash}X@{}}
\toprule
Campo & Contenido \\
\midrule
\texttt{tiempo\_segundos}, \texttt{utilidad\_total} &
Tiempo del algoritmo (sin la fuerza bruta) y
$\sum_t\bigl[f(S_t)+\sum_i u_i\bigr]$ \\
\addlinespace
\texttt{rondas[].seleccion}, \texttt{f\_S},
\texttt{modo\_seleccion}, \texttt{f\_por\_modo} &
Nodos elegidos, su utilidad, el modo ganador y el valor de ambos modos \\
\addlinespace
\texttt{rondas[].traza\_seleccion} &
Por paso: nodo elegido, ganancia marginal, costo y ganancia de
\textbf{todos} los candidatos \\
\addlinespace
\texttt{rondas[].reparto}, \texttt{utilidad\_reparto} &
Unidades por nodo y utilidad $\sum_i u_i(x_i)$ \\
\addlinespace
\texttt{rondas[].traza\_reparto} &
Por unidad: nodo que la recibe y su ganancia marginal \\
\addlinespace
\texttt{rondas[].presupuesto\_restante\_por\_nodo} &
$P_i$ menos lo gastado hasta esa ronda \\
\addlinespace
\texttt{rondas[].exacto} &
Óptimo por fuerza bruta, razón voraz/óptimo y la cota $1-1/e$ \\
\addlinespace
\texttt{rondas[].verificacion} &
Cada restricción con OK/FALLA \\
\bottomrule
\end{tabularx}
\caption{Campos del archivo de resultados.}
\end{table}

\texttt{resultados/<caso>\_traza.txt} contiene lo mismo en texto legible.

%=====================================================================
\section{Pasos de ejecución}
%=====================================================================

Requisitos: Python 3.8 o superior (probado en 3.14). No hay que instalar nada
(\texttt{requirements.txt} no tiene dependencias).

\begin{lstlisting}
git clone <repositorio> && cd federated-learning-optimization
python3 src/red.py data/N15.json      # un caso
python3 src/red.py caso1.json caso2.json -o salida   # varios casos, otra carpeta
python3 src/red.py caso.json --sin-exacto            # sin fuerza bruta (instancias grandes)
python3 tests/test_red.py                            # pruebas
\end{lstlisting}

Cada ejecución imprime la traza y escribe los dos archivos por caso en la
carpeta de salida (por defecto \texttt{resultados/}). El resultado no depende
del directorio desde el que se lance, y dos ejecuciones producen trazas
idénticas.

%=====================================================================
\section{Ejemplo trabajado con el caso de la docente}
%=====================================================================

Datos de la ronda 1 (los valores salen de \texttt{resultados/}).

\paragraph{Selección.}

\begin{table}[htbp]
\centering
\small
\renewcommand{\arraystretch}{1.2}
\begin{tabularx}{\textwidth}{@{}l l r r >{\raggedright\arraybackslash}X@{}}
\toprule
Paso & Elegido & Ganancia marginal & Costo acumulado & Observación \\
\midrule
1 & N11 & 103,490 & 4  & Mayor influencia sobre los demás (pesos de hasta
0,8), pese a tener el menor coeficiente propio (7,6) \\
2 & N14 & 24,278  & 7  & \\
3 & N09 & 9,933   & 10 & \\
4 & N15 & 5,531   & 16 & \\
5 & N13 & 2,348   & 23 & Se alcanza $k=5$; costo $23\le24$ \\
\bottomrule
\end{tabularx}
\caption{Traza de selección de la ronda 1.}
\end{table}

Las ganancias decrecen paso a paso, como corresponde a una función submodular.
Comprobación manual del paso 1:
\[
  f(\{\text{N11}\})
  = a_{\text{N11}}+\sum_{j\ne\text{N11}} a_j\,w_{\text{N11},j}
  =103{,}49.
\]
El voraz por razón ganancia/costo da $f=145{,}106$ y el voraz por ganancia da
$145{,}580$, así que se conserva este último. La fuerza bruta (4\,944
subconjuntos) encuentra el mismo conjunto óptimo, con $f^*=145{,}580$ y razón
1,0.

\paragraph{Reparto de $B=30$.} La primera unidad va a N14 con ganancia
$10{,}1\ln2=7{,}0008$ (la mayor entre los cinco), la segunda a N09 con
$9{,}9\ln2=6{,}8622$, y así sucesivamente. Resultado: $N11=5$, $N14=7$,
$N09=7$, $N15=6$, $N13=5$ (suma 30), con utilidad 88,564.

\paragraph{Efecto del presupuesto total.} Tras la ronda 1, N09 y N14 conservan
11 unidades de 18, y tras la ronda 2 conservan 4. En la ronda 3 solo pueden
recibir 4, así que el reparto cambia a $N11=8$, $N14=4$, $N09=4$, $N15=6$,
$N13=8$ (utilidad 85,611, menor que en las rondas anteriores porque los topes
obligan a mover unidades a nodos de menor ganancia). Al terminar, los cinco
nodos elegidos quedan con 0 de presupuesto.

\textbf{Utilidad total del caso:} 699,477 en 0,001 s.

%=====================================================================
\section{Validación realizada}
%=====================================================================

\begin{enumerate}[leftmargin=*,label=\arabic*)]
  \item \textbf{Verificaciones en cada ejecución} (salen en la traza):
  $|S|\le k$, $\sum c_i\le C$, $\sum x_i\le B$, $x_i\le$ presupuesto
  restante y, como referencia,
  $f(S_{\text{voraz}})\ge(1-1/e)\,f(S_{\text{opt}})$. Todas dan OK en las tres
  rondas.
  \item \textbf{Pruebas automáticas} (\texttt{tests/test\_red.py}):
  \begin{itemize}[leftmargin=*,label=$\bullet$]
    \item el reparto voraz coincide con la enumeración exhaustiva de todos los
    repartos de 8 unidades entre 3 nodos;
    \item $f$ es monótona y submodular en un ejemplo concreto;
    \item se cumplen las restricciones en las 3 rondas y ningún nodo supera su
    presupuesto total de 18.
  \end{itemize}
  \item \textbf{Comparación con el óptimo exacto:} razón voraz/óptimo
  $=1{,}0$ en las tres rondas. Con costos esto es una observación sobre esta
  instancia, no una garantía.
  \item \textbf{Reproducibilidad:} ejecutamos desde otro directorio y con otra
  carpeta de salida, y la traza fue la misma.
\end{enumerate}

\newpage
%=====================================================================
\section{Limitaciones y trabajo futuro}
%=====================================================================

\begin{table}[htbp]
\centering
\small
\renewcommand{\arraystretch}{1.2}
\begin{tabularx}{\textwidth}{@{}>{\raggedright\arraybackslash}p{4.8cm}>{\raggedright\arraybackslash}X>{\raggedright\arraybackslash}p{3.4cm}@{}}
\toprule
Limitación & Consecuencia & Etapa donde se aborda \\
\midrule
$u_i$ y $f$ son elección nuestra (sección~4) &
Los valores numéricos dependen de ella, la corrección del algoritmo no &
Confirmar con la docente \\
\addlinespace
Con costos distintos no hay garantía $(1-1/e)$ &
La razón al óptimo es solo referencial; se necesita el voraz con mejor elemento
individual o enumeración de tripletes &
Algoritmo aproximado (semanas 7-8) \\
\addlinespace
Sin fallos de nodos &
No se modela la variante de rondas con fallos &
Variantes \\
\addlinespace
El reparto ignora las externalidades de red &
Con utilidades dependientes de vecinos, el voraz del reparto puede dejar de ser
óptimo &
Variantes \\
\addlinespace
Selección $O(k\cdot N\cdot m)=O(k\cdot N^2)$ con $f$ explícita &
La $T(n)$ experimental se compara con esta expresión &
Experimentación (semanas 9-10) \\
\addlinespace
Fuerza bruta solo para $N$ pequeño &
No sirve como referencia con instancias grandes &
Algoritmo exacto (semanas 5-6) \\
\addlinespace
Selección miope entre rondas &
No se planifica el uso del presupuesto $P_i$ a lo largo de las $T$ rondas: el
voraz elige igual todas las rondas &
Variantes \\
\addlinespace
Validación de entrada mínima &
Un JSON mal formado produce una excepción de Python &
Sin etapa asignada \\
\addlinespace
Una sola medición de tiempo &
Es válida como evidencia de ejecución, pero es ruidosa &
Experimentación \\
\bottomrule
\end{tabularx}
\caption{Limitaciones de la solución base y su tratamiento posterior.}
\end{table}

\newpage
%=====================================================================
\section{Referencias}
%=====================================================================

\begin{itemize}[leftmargin=*,label=$\bullet$]
  \item B. McMahan, E. Moore, D. Ramage, S. Hampson y B. Agüera y Arcas.
  \emph{Communication-Efficient Learning of Deep Networks from Decentralized
  Data.} AISTATS, 2017.
  \item C. Dwork y A. Roth. \emph{The Algorithmic Foundations of Differential
  Privacy.} Foundations and Trends in Theoretical Computer Science 9(3--4),
  2014.
  \item T. Ibaraki y N. Katoh. \emph{Resource Allocation Problems: Algorithmic
  Approaches.} MIT Press, 1988.
  \item G. L. Nemhauser, L. A. Wolsey y M. L. Fisher. \emph{An analysis of
  approximations for maximizing submodular set functions-I.} Mathematical
  Programming 14, 1978.
  \item U. Feige. \emph{A threshold of $\ln n$ for approximating set cover.}
  Journal of the ACM 45(4), 1998.
  \item S. Khuller, A. Moss y J. Naor. \emph{The budgeted maximum coverage
  problem.} Information Processing Letters 70(1), 1999.
  \item M. Sviridenko. \emph{A note on maximizing a submodular set function
  subject to a knapsack constraint.} Operations Research Letters 32(1), 2004.
  \item R. M. Karp. \emph{Reducibility among combinatorial problems.} En
  \emph{Complexity of Computer Computations}, 1972.
  \item T. H. Cormen, C. E. Leiserson, R. L. Rivest y C. Stein.
  \emph{Introduction to Algorithms} (montículos y algoritmos voraces).
\end{itemize}

\end{document}
